\begin{lemma}[Kempe switch and two-edge closure construction]
Let \(C\) be a \(\noderef{TchangeCounterexample}\) such that its first
color is missing at \(b_3\), and suppose
\(b_1\prec_T b_3\) and \(b_2\prec_T b_3\).  Then there exists a
\(\noderef{TchangeKempeClosureData}\) for \(C\).
\end{lemma}
\begin{proof}
Let \(\sigma\) be obtained by interchanging the two colors \(a,b\) on
the full component containing the third exit path.  The third remote end
misses one of these colors by \(\noderef{TreeExitPath}\),
\(a\) is missing at \(b_3\), and the exact
boundary edge \(f_3\) is colored \(b\).  Hence the third path is the full
component between its two ends: properness and missingness make both
ends degree one, \(\noderef{TwoColorIncidenceDegreeAtMostTwo}\) bounds
all degrees by two, and the spanning path from
\(\noderef{TwoColorComponentSpanningPath}\) cannot have an additional
branch.  The switch therefore makes \(b\) missing
at \(b_3\), by the affected-color clause of
\(\noderef{TwoColorComponentSwapMissingBehavior}\).  Properness follows from
\(\noderef{TwoColorComponentSwapProper}\), and the swap only exchanges
two colored values in the definition
\(\noderef{TwoColorComponentSwap}\), so the unique uncolored edge remains
\(e_0\).

The other two full components are disjoint from the third by
\(\noderef{TwoColorComponentsDisjoint}\): an
intersection would identify the components and put their distinct tree
ends into the component whose first contact with \(T\) is \(b_3\).
Thus their path edges and endpoint missing statuses are fixed.  The third
path remains an exit path with its two colors exchanged.  In all three
cases one obtains walks in the new two-color incidence graph with exactly
the old support and length.

Let \(S=T|b_3\).  The two strict order hypotheses ensure that this prefix
contains \(b_1,b_2\), and it ends at \(b_3\).  The component switch meets
\(T\) only at \(b_3\); its only edge incident with that vertex is the
boundary edge \(f_3\), not a displayed tree edge.  It consequently fixes
every displayed edge and every missing-color witness needed before the
last vertex of \(S\).  Directly from \(\noderef{TashkinovTree}\),
\(S\) is Tashkinov under \(\sigma\).

For \(i=1,2\), let \(x_i\) be the outside endpoint of \(f_i\).  Append
first \((f_1,x_1)\) and then \((f_2,x_2)\).  The \(x_i\) are outside the
old tree by the exact endpoint descriptions in
\(\noderef{BoundaryEdges}\) and \(\noderef{Incident}\), together with
the exit-path condition.  They are distinct: otherwise the
two distinct \(b\)-colored edges would be adjacent by
\(\noderef{AdjacentEdges}\), contradicting
\(\noderef{IsProperEdgeColoring}\).  The edges \(f_1,f_2\) are not
displayed in \(S\), attach at
the earlier vertices \(b_1,b_2\), and have color \(b\), which is missing
at \(b_3\) before and after the first addition.  Hence these are two
successive \(\noderef{TashkinovAugmentation}\) steps.  They produce trees
\(S_1,S_2\), and
\(\noderef{TashkinovAugmentationSequencePreservesTashkinovTree}\) makes
\(S_2\) Tashkinov.

Choose a closure \(U\) of \(S_2\) by
\(\noderef{TashkinovClosureExists}\).  The
\(\noderef{TashkinovClosure}\) definition and another
application of augmentation preservation make \(U\) a closed Tashkinov
tree.  Properness and the one-uncolored-edge condition established above,
together with the coloring-independent clauses of the original
\(\noderef{KTriple}\), give a \(k\)-triple for \(\sigma\).  Therefore
\(\noderef{TashTree}\) says that \(V(U)\) is elementary.  The entrywise
preservation in \(\noderef{TreeSequenceInitialSegment}\) puts all three
\(b_i\) in \(U\), and the first augmentation puts
both ends of \(f_1\) in \(U\).

It remains to prove that no remote end \(a_i\) lies in \(U\).  If one did,
then it misses \(a\) or \(b\), while the distinct vertex \(b_3\) misses
\(b\).  \(\noderef{ColoredVertexSetElementary}\) forces that remote end
to miss \(a\).  Thus both colors are missing somewhere on \(U\).
\(\noderef{ColoredVertexSetClosed}\) forbids either color
on a boundary edge.  Starting at each \(b_j\) and following the
alternating path toward \(a_j\), a first departure from \(U\) would use
exactly such a boundary edge.  Hence all three paths, including all three
distinct remote ends, lie in \(U\).  Elementariness then forces all three
remote ends to miss \(a\), an immediate contradiction.  This proves the
three outside assertions and completes every field of the required
Kempe-closure datum.
\end{proof}
